\section*{Bab \ref{ch:g10:chemical-species} --- Spesies Kimia, Alami dan Sintetis}

\begin{solution}{exo:g10:chemical-species:1}
Alami: vanilin dari polong, kafein dari biji kopi. Sintetis: aspirin,
vanilin pabrik.
\end{solution}

\begin{solution}{exo:g10:chemical-species:2}
Keduanya \omterm{def:g6:pure-substances-and-mixtures:species}{spesies kimia} yang sama, dan sebuah spesies memiliki sifat yang
sama apa pun asalnya.
\end{solution}

\begin{solution}{exo:g10:chemical-species:3}
$R_f = 2.4/6.0 = 0.40$.
\end{solution}

\begin{solution}{exo:g10:chemical-species:4}
Kondensor mendinginkan uap dan mengubahnya kembali menjadi cairan. Jika
masuk dari ujung bawah, air memenuhi seluruh selubungnya (dan bertemu
uap yang paling panas paling akhir, sehingga pendinginannya paling
baik).
\end{solution}

\begin{solution}{exo:g10:chemical-species:5}
Melarutkan spesies itu jauh lebih baik daripada air; \omterm{def:g6:solutions-and-solubility:miscible}{tidak dapat
bercampur} dengan air; sesedikit mungkin berbahaya dan mudah dihilangkan.
\end{solution}

\begin{solution}{exo:g10:chemical-species:6}
Etanol \omterm{def:g6:solutions-and-solubility:miscible}{dapat bercampur} dengan air: tidak terbentuk lapisan, tidak ada
yang dapat dipisahkan. Sikloheksana \omterm{def:g6:solutions-and-solubility:miscible}{tidak dapat bercampur} dengan air dan
melarutkan linalool dengan baik.
\end{solution}

\begin{solution}{exo:g10:chemical-species:7}
Air berada di atas: \qty{1}{mL} diklorometana beratnya \qty{1.33}{g},
lebih berat daripada \qty{1}{mL} air, sehingga diklorometana tenggelam.
\end{solution}

\begin{solution}{exo:g10:chemical-species:8}
Serbuk itu meleleh terlalu rendah (aspirin: \qty{135}{\celsius}) dan
dalam suatu rentang: sampelnya tidak murni, atau bukan aspirin.
\end{solution}

\begin{solution}{exo:g10:chemical-species:9}
Linalool dan linalil etanoat (kedua bercaknya setinggi kedua
pembanding). Linalil etanoat naik lebih tinggi.
\end{solution}

\begin{solution}{exo:g10:chemical-species:10}
\omterm{def:g10:chemical-species:distillation}{Hidrodistilasi}; menambahkan sikloheksana ke distilat dan mengocoknya;
memisahkan lapisan-lapisan; \omterm{def:g3:separating-mixtures:evaporate}{menguapkan} sikloheksana.
\end{solution}

\begin{solution}{exo:g10:chemical-species:11}
Pensil tidak \omterm{def:g2:mixing-and-dissolving:dissolve}{larut} dalam \omterm{def:g10:chemical-species:tlc}{eluen}; di atas \omterm{def:g10:chemical-species:tlc}{eluen}, totolan-totolan itu
terbawa naik melalui pelat alih-alih \omterm{def:g2:mixing-and-dissolving:dissolve}{larut} ke dalam bejana.
\end{solution}

\begin{solution}{exo:g10:chemical-species:12}
Sikloheksana mendidih pada \qty{81}{\celsius}, linalool pada
\qty{198}{\celsius}: pemanasan pelan mengusir sikloheksana sementara
linaloolnya tertinggal.
\end{solution}

\begin{solution}{exo:g10:chemical-species:13}
Sampel itu kemungkinan spesies tersebut ($R_f$ sama dengan \omterm{def:g10:chemical-species:natural}{spesies
alami} pada pelat yang sama) dan kemungkinan murni (satu bercak, titik
leleh tajam).
\end{solution}

\begin{solution}{exo:g10:chemical-species:14}
$0.40 \times 5.5 = \qty{2.2}{cm}$. Bercak lainnya pada $0.75 \times 5.5 =
\qty{4.1}{cm}$ (sampai mm terdekat), jadi sekitar \qty{1.9}{cm} di atas
A.
\end{solution}

\begin{solution}{exo:g10:chemical-species:15}
\qty{200}{mL} dapat melarutkan $1.6 \times 0.200 = \qty{0.32}{g}$:
dengan \qty{0.20}{g} distilat itu belum jenuh; masih dapat melarutkan
\qty{0.12}{g} lagi. \qty{0.20}{g} yang \omterm{def:g2:mixing-and-dissolving:dissolve}{larut} itu akan hilang bersama
airnya; sikloheksana, yang melarutkan linalool jauh lebih baik, menarik
kembali sebagian besarnya.
\end{solution}

\begin{solution}{pb:g10:chemical-species:1}
\textbf{1.} Uapnya membawa spesies-spesies harum keluar dari bunga dan
masuk ke kondensor.

\textbf{2.} Kondensor mendinginkan uap kembali menjadi cairan; air
dingin masuk dari ujung bawah.

\textbf{3.} \omterm{def:g6:solutions-and-solubility:miscible}{Tidak dapat bercampur}; massa jenisnya lebih kecil daripada
air (minyak itu terapung).

\textbf{4.} Lapisannya sangat tipis dan sebagian minyaknya \omterm{def:g2:mixing-and-dissolving:dissolve}{larut} dalam
air atau tersebar di dalamnya: \omterm{def:g10:chemical-species:extraction}{ekstraksi} memperolehnya kembali.

\textbf{5.} Sikloheksana melarutkan linalool dengan baik, \omterm{def:g6:solutions-and-solubility:miscible}{tidak dapat
bercampur} dengan air, dan mendidih pada suhu rendah
(\qty{81}{\celsius}), sehingga mudah dihilangkan.

\textbf{6.} Tidak: etanol bercampur dengan air dan tidak terbentuk
lapisan.

\textbf{7.} Di atas (\qty{0.78}{g} per mL, kurang daripada air). Dengan
diklorometana (\qty{1.33}{g} per mL) lapisan \omterm{def:g6:solutions-and-solubility:solute}{pelarutnya} akan berada di
bawah.

\textbf{8.} Tidak ada nyala api di dekatnya (mudah menyala); bekerja di
dalam lemari asam dengan sarung tangan, dan mengumpulkan limbahnya
alih-alih membuangnya ke bak cuci (beracun bagi kehidupan akuatik).

\textbf{9.} Sikloheksana mendidih pada \qty{81}{\celsius}, linalool pada
\qty{198}{\celsius}: pemanasan pelan hanya menghilangkan \omterm{def:g6:solutions-and-solubility:solute}{pelarutnya}.

\textbf{10.} $R_f(\text{L}) = 2.4/6.0 = 0.40$; $R_f(\text{A}) = 4.5/6.0
= 0.75$.

\textbf{11.} Linalool dan linalil etanoat.

\textbf{12.} Bukan: minyak itu mengandung sedikitnya dua spesies; minyak
itu \omterm{def:g6:pure-substances-and-mixtures:mixture}{campuran}.

\textbf{13.} Agar semua bercak dikembangkan dalam kondisi yang sama:
hanya dengan begitu ketinggiannya dapat dibandingkan.

\textbf{14.} $R_f$-nya sama dengan \omterm{def:g10:chemical-species:retention-factor}{faktor retensi} linalool dalam kondisi
yang sama: sampel itu kemungkinan besar linalool, identik dengan \omterm{def:g10:chemical-species:natural}{spesies
alaminya}.

\textbf{15.} $1.0 \times 0.90 = \qty{0.90}{g}$.

\textbf{16.} $0.90/100 = \qty{0.90}{\percent}$.

\textbf{17.} $100/0.009 \approx 11\,000$\,g, sekitar \qty{11}{kg}
bunga.

\textbf{18.} Minyak itu diekstraksi dari tumbuhan, tanpa \omterm{def:g5:irreversible-changes:chemical-change}{perubahan kimia}.
Tetapi linaloolnya adalah spesies yang sama dengan linalool sintetis:
rumus yang sama, sifat yang sama.

\textbf{19.} \qty{0.90}{g} minyak atsiri dari \qty{100}{g} bunga.
\end{solution}
